\section{Exceptional Control Flow and Error Routing}

Not every execution route fails locally. Some operations fail in a way that cannot be handled by simply moving to the next statement or checking a returned value. Python uses exceptions for this kind of route change. An exception interrupts the ordinary path and sends execution toward a matching handler, possibly outside the current block or even outside the current function call.

This section studies exceptions as control-flow structures, not merely as error messages. A raised exception is a non-local jump through the active execution graph. The exception object carries a runtime type, the handler selection mechanism matches that type against exception classes, and the traceback records the route through active frames.

The goal is to make exception handling precise rather than magical: \texttt{try} marks a protected region, \texttt{except} catches selected failure routes, \texttt{else} isolates success-only logic, \texttt{finally} guarantees cleanup, and \texttt{raise} deliberately creates or continues an exceptional edge.